\documentclass[11pt]{article}
\usepackage[margin=2.5cm]{geometry}
\usepackage{amsmath,amssymb}
\usepackage{graphicx}
\graphicspath{{./}{../figs/}{figs/}}
\usepackage{booktabs}
\usepackage{url}
\usepackage[numbers]{natbib}
\setlength{\emergencystretch}{1em}
\newcommand{\rhosp}{\ensuremath{\rho}}
\renewcommand{\thesection}{S\arabic{section}}
\renewcommand{\thetable}{S\arabic{table}}
\renewcommand{\thefigure}{S\arabic{figure}}
\renewcommand{\theequation}{S\arabic{equation}}
\title{Supplementary Material for\\ ``Order without Magnitude: Single-Unit Fine-Tuning Can Make
Vision--Language Models Rank Motion but Ignore the Requested Unit''}
\author{Jun Ji, Bowen Tan, Yizhou Zhao, Yi Li, Xiaolei Zhang, Yi Sui}
\date{}
\begin{document}
\maketitle
\noindent This supplement collects the protocol details, provenance records and
secondary analyses behind the main text. S1 documents training and provenance for
every checkpoint; S2 the prompts, the text probe and the paired analyses; S3 the benchmark
construction and references; S4 the cross-question consistency check, the matched-label
comparison and the ordering results by viewpoint; S5 the per-backbone, read-then-convert,
probe and zero-shot results and the per-item conversion ratios. Analyses of the earlier version of this study that the main
text no longer uses are archived with the released records (end of S4). References to the
main text are given in words.

\section{Training details and provenance}
\label{app:adapters}

\subsection{Original CourtDyn-native: trained within this study}

CourtDyn-native starts from the public Qwen3.5-4B checkpoint with no adapter
initialization. It uses NF4 double quantization with bf16 compute, LoRA
($r=16$, $\alpha=32$, dropout $0.05$) on the language-tower
q, k, v, o, gate, up and down projections, and a frozen vision tower. Image processing
is capped at $200{,}704$ pixel-equivalent. Training uses CourtDyn v1 speed/path
answers from four side clips and Q4\_top, 1,400 items in total, for one epoch
(88 optimizer steps; batch 1, accumulation 16; learning rate $10^{-4}$).
The initialized backbone, task data, label generator and deterministic scorer
are specified here; no static-task training pool or transferred adapter is
required for this model.

Q1\_top and Q2\_top contribute no camera-view frame or QA item to this
training list. However, Q1\_side and Q1\_top depict the same opening event.
The primary Q2 case is temporally disjoint from the Q2 training clip, within
the same game. A corrected builder excludes Q1\_side and prepares a separate
1,120-item pool with no identified test-event overlap. The original results
retain the original 1,400-item checkpoint and its documented split.

\subsection{Matched event-holdout training}
\label{app:matched_label_training}

The v1-trained and v3-trained adapters both completed one epoch on the
1,120-item pool (70 optimizer steps). Ordered questions, images and item
identities match, while the supervisory answers follow v1 or v3. Both runs
start from the same public backbone with no adapter initialization and set
seed 42 before model loading. They use NF4 loading with bf16 compute, a
frozen vision tower and language-tower LoRA ($r=16$, $\alpha=32$, dropout
$0.05$), batch 1, accumulation 16, learning rate $10^{-4}$ and the same
$200{,}704$ pixel budget. No invalid training items are dropped.
The completion records identify the final weights, training inputs and
protocols. The tables of this section are seed 42; seeds 43 and 44 of both pools
were trained later with identical settings, and the main text reports the three
seeds together. Comparing either new adapter with the original 1,400-item,
88-step checkpoint does not isolate the effect of removing event overlap.

The evaluation protocol applies the explicit-coordinate metre, centimetre
and pixel prompts with the height sentence retained to Q1 and Q2. It also
pairs metre prompts with four copies of the first frame and includes the
same 12 text-only arithmetic probes on Q2. It does not repeat the no-height
factorial or the legacy ruler intervention for these new adapters.
Each temporal comparison uses the matching new full predictions from
the same checkpoint; historical full outputs are not substituted. The
completed new-model evaluation contains 4,480 visual answers, of which
4,472 strictly parse, plus 12 text-only answers per checkpoint.
Saved adapter configurations differ only in the serialization order of
the same seven distinct target-module names. Installed PEFT converts this
list to a set and matches modules by membership or suffix. A separately
audited compatibility wrapper sorts deep copies for the configuration
comparison only, retains every other check, restores the originals and
preserves the original order-sensitive failure report. It changes no
checkpoint bytes, predictions or numerical scoring.

\subsection{Completed matched-label comparisons}
\label{app:matched_label_results}

The new-adapter unit results (Table~\ref{tab:matched_label_units}) use the
same strict parser and exact references as the original explicit-coordinate
control. The nominal pixel tolerance factors are $K=27.061$ for Q1 and
$K=27.498$ for Q2; these scale $T$ and $\epsilon$, not the itemwise reference
ratio. All full-frame unit comparisons have 140 common positive metre
denominators. The first-frame results (Table~\ref{tab:matched_label_temporal})
keep all-item scores separate from common-parsed correlations. Q1 v3 speed
has five invalid control answers; its 135 remaining values vary, so its
correlation is defined. Q2 v3 speed instead has three invalid answers and
137 identical valid predictions. Neither case is missing inference.

\begin{table}[t]
\centering\small
\setlength{\tabcolsep}{4pt}
\caption{New adapters under retained-height unit prompts. Entries for cm
and px give all-item T-MRA / Spearman correlation; all cells have 140 valid
answers. $R_{px/m}$ is the median paired projected-pixel/metre prediction
ratio. Its exact reference-ratio medians are 27.10/27.06 (Q1 speed/path)
and 27.49/27.50 (Q2). These differ from the exact cm/m factor of 100.}
\label{tab:matched_label_units}
\begin{tabular}{llrrr}
\toprule
Train/event & Family & cm T/$\rho$ & px T/$\rho$ & $R_{px/m}$ \\
\midrule
v1 / Q1 & speed & 15.00/0.406 & 16.79/0.394 & 1.56 \\
v1 / Q1 & path & 11.36/0.568 & 15.64/0.410 & 3.47 \\
v1 / Q2 & speed & 21.36/0.619 & 23.43/0.572 & 1.67 \\
v1 / Q2 & path & 20.43/0.549 & 23.43/0.630 & 1.55 \\
v3 / Q1 & speed & 14.86/0.360 & 16.43/0.619 & 1.57 \\
v3 / Q1 & path & 11.07/0.753 & 12.21/0.224 & 1.00 \\
v3 / Q2 & speed & 21.57/0.540 & 23.07/0.634 & 1.83 \\
v3 / Q2 & path & 20.21/0.737 & 21.71/0.429 & 1.45 \\
\bottomrule
\end{tabular}
\end{table}

\begin{table}[t]
\centering\small
\setlength{\tabcolsep}{3pt}
\caption{New-adapter first-frame comparison. Both T-MRA columns use all
140 observed items with invalid answers scoring zero. $N$ is the common
parsed cohort for both correlations; it is not the T-MRA denominator.
``u'' denotes an undefined correlation: all 137 parsed v3 Q2 speed control
answers equal 0.1. Full denotes the matching new metre/height arm.}
\label{tab:matched_label_temporal}
\begin{tabular}{llrrrrr}
\toprule
 & & & \multicolumn{2}{c}{T-MRA} & \multicolumn{2}{c}{$\rho$} \\
\cmidrule(lr){4-5}\cmidrule(lr){6-7}
Train/event & Family & $N$ & Full & First & Full & First \\
\midrule
v1 / Q1 & speed & 140 & 54.14 & 32.00 & 0.600 & 0.057 \\
v1 / Q1 & path & 140 & 47.21 & 29.29 & 0.512 & -0.169 \\
v1 / Q2 & speed & 140 & 72.07 & 43.79 & 0.715 & 0.202 \\
v1 / Q2 & path & 140 & 74.21 & 44.00 & 0.694 & 0.038 \\
v3 / Q1 & speed & 135 & 51.93 & 18.86 & 0.658 & 0.076 \\
v3 / Q1 & path & 140 & 47.50 & 18.86 & 0.778 & -0.241 \\
v3 / Q2 & speed & 137 & 67.64 & 26.57 & 0.759 & u \\
v3 / Q2 & path & 140 & 64.14 & 29.36 & 0.773 & 0.036 \\
\bottomrule
\end{tabular}
\end{table}

The 12 text-only probes are identical across checkpoints, not 36 independent
stimuli. The v1-trained adapter gives correct cm answers for 0.2, 0.5 and
2.5 metres, and correct mm answers for all six values. Its other cm answers
are ten times the metre input instead of one hundred times. The v3-trained
adapter gives only the 2.5-metre cm answer correctly; mm answers are correct
for 0.5, 1 and 8 metres. Each of its eight errors is ten times smaller than
the requested target. These are output patterns, not identified internal
conversion algorithms.

\subsection{Same-prediction dual-reference scoring}
\label{app:paired_reference_scores}

We reconstruct exact v1 and v3 answers from the same annotation inputs for
all 560 held-out speed/path items, preserving the original smoothing,
time intervals and reference definitions. The metre predictions from both
new adapters and both frame conditions are each scored against both
references: 32 exact cells, without new inference. A separate 32-cell
one-decimal sensitivity retains the historical rounding convention.
Both exact metre references use the same family $T$ and $\epsilon$ from
Table~\ref{tab:families}; changing reference convention is not unit
rescaling. For the common strictly parsed set $P$, we additionally report
\begin{equation}
\label{eq:paired_reference_mae}
\mathrm{MAE}_{P}=\frac{1}{|P|}\sum_{i\in P}|\hat y_i-y_i|.
\end{equation}
Invalid answers are excluded from MAE, never assigned zero error; its
denominator and unit are reported separately. All-item T-MRA still gives
invalid answers zero score. Within-model reference comparisons use identical
predictions. Paired model and temporal correlations use common parsed
cohorts. The full-frame MAE comparison uses the same 140 parsed items per
arm; supplementary per-arm MAE retains each arm's own parsed cohort.
Constant predictions retain undefined correlations under either reference.

Table~\ref{tab:paired_reference_scores} gives the eight full-frame rows;
the aggregate supplement retains both frame conditions and the rounding
sensitivity. Both models are closer in MAE to v3 than v1, including the
v1-trained adapter. The v3-trained adapter nevertheless has higher MAE
than the v1-trained adapter in every event/family under either reference.
Thus reference choice strongly affects the absolute scores, and matching
training labels to the homography does not guarantee better magnitude fit
in this seed-42 comparison (dual-reference scoring was not repeated for seeds 43
and 44). The v3 reference remains subject to the
calibration and foot-point limitations of Section~\ref{sec:gt}.

\begin{table}[t]
\centering\small
\setlength{\tabcolsep}{4pt}
\caption{Same full-frame metre predictions scored against both exact
reference conventions. Each row uses the same 140 observed and valid
answers for T-MRA and MAE. MAE is in m/s for speed and m for path; it is
not a percentage. Training labels identify the model, while column groups
identify the scoring reference.}
\label{tab:paired_reference_scores}
\begin{tabular}{llrrrr}
\toprule
 & & \multicolumn{2}{c}{exact v1} & \multicolumn{2}{c}{exact v3} \\
\cmidrule(lr){3-4}\cmidrule(lr){5-6}
Train/event & Family & T-MRA & MAE & T-MRA & MAE \\
\midrule
v1 / Q1 & speed & 21.57 & 2.985 & 54.14 & 0.963 \\
v1 / Q1 & path & 23.57 & 6.663 & 47.21 & 2.096 \\
v1 / Q2 & speed & 27.71 & 1.812 & 72.07 & 0.586 \\
v1 / Q2 & path & 32.50 & 3.864 & 74.21 & 1.210 \\
v3 / Q1 & speed & 17.71 & 3.037 & 51.93 & 0.972 \\
v3 / Q1 & path & 15.36 & 7.094 & 47.50 & 2.220 \\
v3 / Q2 & speed & 20.93 & 1.887 & 67.64 & 0.608 \\
v3 / Q2 & path & 21.21 & 4.285 & 64.14 & 1.395 \\
\bottomrule
\end{tabular}
\end{table}


\subsection{Adapters for the seed, control and backbone experiments}
\label{app:revision_adapters}

\textbf{Seeds.} Seeds 43 and 44 of the CourtDyn-native, v1 and v3 pools reuse
the pools, order and hyperparameters of seed 42; only the seed set before model
loading differs.

\textbf{Mixed-unit adapters.} The v3 pool was rebuilt with 560 of its 1,120 items
asking for centimetres and carrying the metre answer multiplied by 100 (exact
decimal arithmetic), balanced within every clip and family and selected by a
content hash; questions, images, order and hyperparameters are otherwise those of
the v3 pool. The build receipt records the SHA-256 of both pools. Seeds 42, 43
and 44 were trained (70 optimizer steps each).

\textbf{Other backbones.} SmolVLM2-2.2B and Qwen2.5-VL-3B were first trained with
the one-epoch recipe above (round 1, 70 steps). In round 2 each was trained on a
four-epoch schedule (280 steps) with a checkpoint every 70 steps, and the epoch was
chosen on the 280-item development set by the rule of the main text: $k^*=4$ for
both, with SmolVLM2 degenerate at the cap. InternVL3-2B was run under the round-2
protocol only ($k^*=2$), with one $448\times448$ image tile per frame and its own
unadapted-backbone cells. LoRA targets the same language-tower projections in every
backbone, with the vision tower frozen; the mixed-unit adapter of each backbone is
taken at that backbone's $k^*$.

\subsection{Additional frozen transfer checkpoints}

Three previously trained checkpoints over the same backbone appear in the
archived analyses of the earlier version of this study and in the released
records. Their names identify different training procedures, not separate
backbone families; their configurations are given here so that those records
can be read without another manuscript.

\textbf{source-SFT.} One epoch on 20k sport-stratified geometry-generated QA
from calibrated racket-sport images: 1,250 steps, batch 1 with accumulation 16,
learning rate $10^{-4}$, cosine schedule and 3\% warmup.

\textbf{GRPO-v3.} Continued from source-SFT on 1,855 static basketball
prompts, with only LoRA parameters updated. Group size 6, temperature 1.0,
top-$p$ 0.95, up to 48 new tokens, learning rate $10^{-5}$, 20 warmup steps,
advantage clipping $\pm2$ and gradient clipping 1.0; no KL penalty. Reward
uses a geometry-based answer scorer. These settings do not formally bound
policy drift.

\textbf{target-SFT-token.} Supervised training from source-SFT on the same
static basketball prompt pool, matching GRPO's 106,514 applied completion-label
tokens over 681 steps (seed 42), with the same adapter, quantization and
optimizer configuration. This is a frozen training-history comparison;
we import no earlier accuracy or causal conclusion into the present analysis.

\subsection{Related preprint and shared material}
\label{app:shared}

These three checkpoints also appear in a preprint, currently under review,
by an overlapping author group on static metric spatial
reasoning~\cite{ji2026calibration}. They are reused here as frozen comparators,
with new dynamics QA and evaluations. Their provenance is disclosed so that
these comparators are not mistaken for models trained on CourtDyn. The main
CourtDyn-native model was instead trained from the public unadapted backbone.

Both studies use Human-M3. Of the 224 unique camera--frame pairs in the four
basketball1/split1 pools evaluated here, 73 also occur in that study's static
evaluation. These are shared source frames, not identical QA or identical
rendered inputs. The additional Human-M3 experiment is therefore external to
TeamTrack, but not an entirely unused image source relative to the other
study. No static-task score or claim from that preprint is used as
evidence for the main result in this paper.

\subsection{Reconstruction boundary}

The original native adapter is listed in the paper2 artifact archive; the three
additional transfer adapters are not included in that archive. Their frozen
CourtDyn outputs support checking the reported comparisons once the licensed
pools are rebuilt, but rerunning their inference requires the exact checkpoint
weights. Reconstructing their static training history is not part of the
native-model workflow. The primary result and the supplemental comparisons
therefore have different reproduction prerequisites, which should not be
conflated.

\section{Prompt and paired-analysis details}
\label{app:paired}

The original speed prompt is:
\begin{quote}\small
The 4 frames are consecutive samples from one basketball clip spanning
2.0 seconds, in chronological order. Assume a typical player on this court is
1.93 m tall. What is the average speed of the player marked with the red box,
in meters per second? Output only the number, one decimal place.
\end{quote}
For path, the question sentence is ``How far does the player marked with the
red box travel along the floor over the whole clip, in meters?'' The noprior
arm deletes the height sentence. The ruler arm replaces it with ``In these
frames, one meter on the floor spans about 27 pixels.'' The pxunit arm deletes
the height sentence and replaces the speed question by ``What is the average
speed of the player marked with the red box across the image, in pixels per
second?'' For path, it asks how far the player moves across the image over the
whole clip, in pixels. None of these original prompts specifies the uploaded
canvas size. We report the original prompts, rather than silently substituting
a corrected prompt for frozen predictions.

\paragraph{New explicit-coordinate controls}
\label{app:explicit_unit_controls}
The explicit prompts of the main text retain the original four frames and
two-second lead, then insert: ``Each supplied frame is 960 pixels wide and
540 pixels high. Pixel answers refer to this original canvas before any
resizing. Physical units describe floor-plane motion; pixels describe
projected image-plane motion.'' The height sentence is either retained
verbatim or deleted. The speed/path request uses metres, centimetres or
pixels (per second for speed), followed by the original one-decimal
number-only instruction. All pairs use the same ordered images and
track/window. Metre targets are exact v3 values; centimetre targets and their
$T,\epsilon$ are multiplied by $100$. Pixel targets are measured on rendered
image trajectories, with $T,\epsilon$ multiplied by nominal $K=27.498$ on Q2.
T-MRA (main text) uses unrounded references and parsed numbers; a valid
answer must be a finite numeric string in its entirety. Invalid parses would
score zero in the all-item score; none occurs here. Spearman uses these exact
references. $R$ is the paired ratio of the main text, with the metre
arm under the same height condition as denominator; all 140 denominators are positive (Table~\ref{tab:explicit_unit_controls}).

The no-image arithmetic probe uses metre inputs
$\{0.2,0.5,1.0,2.5,4.0,8.0\}$, each requested in centimetres and millimetres
with the factor $100$ or $1000$ supplied. The checkpoint, greedy decoding,
4-bit loading and final-answer suffix match the visual run. Correctness
requires a parsed answer with absolute error less than $0.05$. For the original
native checkpoint both correct answers are millimetre conversions ($0.5\to500$,
$2.5\to2500$); the remaining millimetre outputs are tenfold or hundredfold too
small. These twelve items are the original probe; they are contained verbatim in
the 392-item probe below, which supersedes them for every conclusion in the
main text.

\begin{table}[t]
\centering\small
\setlength{\tabcolsep}{4pt}
\caption{New explicit-coordinate controls, original native checkpoint, Q2.
Each family/arm has 140 observed and parsed answers. T-MRA is an all-item
tolerance score; $R$ compares each item with its metre prediction under the
same height condition, with metre rows serving as the reference.
The listed units apply to path; speed uses the corresponding unit per second.}
\label{tab:explicit_unit_controls}
\begin{tabular}{llrrrrrr}
\toprule
 & & \multicolumn{3}{c}{speed} & \multicolumn{3}{c}{path} \\
\cmidrule(lr){3-5}\cmidrule(lr){6-8}
Unit & Height & T-MRA & $\rho$ & $R$ & T-MRA & $\rho$ & $R$ \\
\midrule
m & yes & 77.9 & 0.699 & --- & 70.9 & 0.683 & --- \\
cm & yes & 21.3 & 0.683 & 1.00 & 20.2 & 0.714 & 1.02 \\
px & yes & 23.1 & 0.639 & 1.60 & 21.6 & 0.735 & 1.80 \\
m & no & 78.6 & 0.699 & --- & 62.4 & 0.642 & --- \\
cm & no & 21.3 & 0.625 & 1.00 & 20.2 & 0.705 & 1.00 \\
px & no & 23.0 & 0.654 & 1.60 & 21.5 & 0.707 & 1.80 \\
\bottomrule
\end{tabular}
\end{table}

\paragraph{Expanded text-only unit probe}
\label{app:unit_probe}
The probe crosses fourteen magnitudes
$\{0.002,\allowbreak 0.007,\allowbreak 0.02,\allowbreak 0.05,\allowbreak 0.2,\allowbreak 0.5,\allowbreak 1.0,\allowbreak 2.5,\allowbreak 4.0,\allowbreak 8.0,\allowbreak 17.5,\allowbreak 40,\allowbreak 125,\allowbreak 400\}$
with fourteen conditions: metres to and from centimetres, millimetres and
kilometres, each with the conversion factor stated or withheld, and metres to
and from pixels at $K=27.498$ with the factor always stated, since $K$ is not
recallable knowledge. Every item appears bare (``A distance is $x$ metres.'')
and embedded in the CourtDyn frame (``The player marked with the red box has a
path length of $x$ metres over the clip.''), giving $392$ items. Values are
written exactly as the published probe writes them (``1.0'', not ``1''), so
its twelve items appear verbatim; the generator refuses to write the file
otherwise. Correctness keeps the published criterion, $|\mathrm{error}|<0.05$
with a one-decimal answer, so the published items keep their published scores.
That criterion was written for multiplications to targets of $20$ or more.
On $146$ items, all divisions, an unconverted answer would also score as
correct (for example $0.02$~m to kilometres, target $0.00002$), so those items
cannot detect the failure the probe exists to find. We report accuracy on the
$246$ remaining diagnostic items; the all-item figures are retained with the
analysis outputs. Wilson intervals are reported per proportion. Adapter and
backbone are compared with an exact McNemar test on the same items. Each
answer is also classified by $\log_{10}(\text{prediction}/\text{target})$ as
correct, unchanged input, reversed direction, an error of a whole power of
ten, or other. Every checkpoint ran as its own process under the visual
protocol with an empty image list.

\paragraph{Two-stage unit answering}
\label{app:two_stage}
Stage one is not rerun: each checkpoint's published Q2 metre-arm answers are
reused, and the raw answer string is inserted verbatim into the second turn
(``The player marked with the red box has a path length of $x$ meters over
the clip. One meter equals 27.498 pixels in this image. Express that path
length in pixels.'', and the corresponding speed and centimetre forms). The
second turn carries no image and uses the visual protocol otherwise. Its
answers are scored with the same T-MRA and strict parser against the target
arm's own references, so direct, two-stage and exact-conversion scores share
one scale. Items with an unparsed metre answer have nothing to convert and
score zero in all three. The predictions and thresholds were written before
any second-turn output existed and are decided mechanically by the analysis
script (labelled P1--P3 in the record, H1--H3 here to avoid confusion with
the protocols of the main text): H1, ratio within $5\%$ of $K$
and T-MRA within $5$ points of the ceiling; H2, ratio not within $5\%$ of
$100$ and T-MRA at least $20$ points below the ceiling; H3, ratio within
$5\%$ of $100$. Two amendments were logged
before the second turn ran: a fix to a file-completeness check, and the
observation that the backbone's metre answers are degenerate, which makes its
T-MRA clause uninformative. The record is local and hashed, not externally
timestamped.

\paragraph{Historical v1 first-frame scoring} The original first-frame
summary scored both arms against v1 rather than v3. For Q2, speed/path
correlations change from $0.74$/$0.75$ to $0.16$/$-0.06$; for Q1 they
change from \mbox{$0.72\to0.03$} and \mbox{$0.69\to-0.05$}. The full-minus-control
range across these four v1 cells, i.e.\ $\Delta\rho$ of $+0.58$ to $+0.81$,
therefore differs slightly from the v3 result in
the main text. Predictions and sample membership are the
same; only the reference-answer convention changes.

\paragraph{Sensitivity analysis} To describe the paired data more
fully, we recomputed ratios and their interquartile ranges on common valid
items. All 16 adapter pixel-unit cells have 140 common parsed predictions and
140 positive full-arm denominators. We also resampled whole 5 s bins of window
start times, including all players in a bin, using 2,000 paired bootstrap
replicates (seed 20260912). Table~\ref{tab:paired-detail} reports the resulting
percentile intervals for the difference between pixel-target rho and v3
metre-target rho. An accompanying 10 s-bin sensitivity is retained with the
analysis outputs. These are exploratory analyses, not preregistered tests:
there are only six or three bins per clip, boundary windows and players can
still link bins, and no equivalence margin was specified. The intervals
therefore do not certify independent-game generalisation or unchanged ordering.

\begin{table}[t]
\centering\footnotesize
\setlength{\tabcolsep}{3pt}
\caption{Pixel-unit paired detail. $R$ is the paired ratio of the main text; brackets
  in its column give the IQR, not a confidence interval. The last column gives
  $\Delta\rho$ and an exploratory 5 s-bin bootstrap 95\% interval.}
\label{tab:paired-detail}
\begin{tabular}{llllll}
\toprule
Clip & Model & Family & $n$ & $R$ [IQR] & $\Delta\rho$ [interval] \\
\midrule
Q1 & source-SFT & speed & 140 & 2.00 [1.40, 2.50] & -0.08 [-0.22, +0.02] \\
Q1 & source-SFT & path & 140 & 1.00 [0.84, 1.37] & -0.03 [-0.13, +0.05] \\
Q1 & GRPO-v3 & speed & 140 & 1.30 [1.11, 1.47] & +0.05 [-0.01, +0.13] \\
Q1 & GRPO-v3 & path & 140 & 1.03 [0.97, 1.18] & -0.02 [-0.07, +0.04] \\
Q1 & target-SFT-token & speed & 140 & 1.30 [1.10, 1.73] & +0.18 [+0.08, +0.30] \\
Q1 & target-SFT-token & path & 140 & 1.27 [1.15, 1.52] & +0.03 [-0.03, +0.18] \\
Q1 & CourtDyn-native & speed & 140 & 1.08 [1.00, 1.40] & -0.03 [-0.10, +0.03] \\
Q1 & CourtDyn-native & path & 140 & 0.62 [0.53, 1.00] & -0.08 [-0.17, +0.07] \\
Q2 & source-SFT & speed & 140 & 1.50 [1.08, 2.25] & +0.05 [-0.06, +0.16] \\
Q2 & source-SFT & path & 140 & 1.00 [1.00, 1.40] & -0.00 [-0.07, +0.06] \\
Q2 & GRPO-v3 & speed & 140 & 1.36 [1.10, 1.73] & -0.03 [-0.09, +0.03] \\
Q2 & GRPO-v3 & path & 140 & 1.00 [0.90, 1.10] & +0.05 [+0.00, +0.12] \\
Q2 & target-SFT-token & speed & 140 & 1.10 [1.00, 1.80] & +0.04 [-0.05, +0.17] \\
Q2 & target-SFT-token & path & 140 & 1.25 [1.00, 1.33] & +0.04 [-0.00, +0.11] \\
Q2 & CourtDyn-native & speed & 140 & 1.11 [1.00, 1.60] & -0.06 [-0.09, -0.01] \\
Q2 & CourtDyn-native & path & 140 & 0.62 [0.53, 1.00] & -0.03 [-0.08, +0.03] \\
\bottomrule
\end{tabular}
\end{table}


\paragraph{Cross-question consistency}
\label{app:consistency}
We pair speed and path questions by window and player track, verifying that
their image identifiers and frame lists agree. Each family contains 140
questions, but the common trajectory sets contain only 100 pairs in Q2 and
121 in Q1. The full, ruler and pixel arms shown use the same common sets.
For these pairs, the unrounded v1, v3 and pixel references all satisfy
$d=\Delta t v$ to numerical precision. The actual duration is
$\Delta t=2.002$ s (60 frame intervals at 29.97002997 fps); the prompt
rounds this to 2.0 s. For positive speed predictions we define
\begin{equation}
 S_i=\frac{\hat d_i}{\Delta t_i\hat v_i},\qquad
 C_i=\mathbf{1}\!\left\{
 |\hat d_i-\Delta t_i\hat v_i|\leq0.05(1+\Delta t_i)\right\}.
\end{equation}
The bound allows rounding both answers to one decimal place in the requested
units. It therefore imposes different physical tolerances in metres and
pixels: unlike $S$, $C$ is not invariant to a common rescaling of both
predictions. We report $C$ only as compatibility with the requested decimal
precision, not as an accuracy threshold, a prespecified equivalence margin,
or a unit-invariant measure of consistency.
All listed pairs (Table~\ref{tab:cross-question}) have finite predictions and positive speed denominators.
Using 2.0 s and its corresponding bound of 0.15 instead of 0.1501 leaves
every listed pair's classification unchanged. A shared unknown coordinate
scale cancels from $S$, so uniform rescaling alone cannot account for its
change across prompts. Different conventions between questions, answer priors
or wording effects remain possible. These are post-hoc descriptive analyses
of dependent observations; Q1 contains training-event overlap.

\begin{table}[t]
\centering\small
\caption{Native speed--path consistency on common trajectories.
Brackets are item-level IQRs. $n_C$ counts compatibility with decimal
rounding within each requested unit, whose physical tolerances differ.
It does not count correct motion estimates; even constants can be consistent.}
\label{tab:cross-question}
\begin{tabular}{llrll}
\toprule
Clip & Arm & Pairs & Median $S$ [IQR] & $n_C$ \\
\midrule
Q2 & full & 100 & 0.999 [0.999, 1.155] & 42/100 \\
Q2 & ruler & 100 & 0.999 [0.749, 0.999] & 53/100 \\
Q2 & pixel & 100 & 0.500 [0.500, 0.659] & 4/100 \\
Q1 & full & 121 & 1.135 [0.999, 1.499] & 21/121 \\
Q1 & ruler & 121 & 0.999 [0.666, 0.999] & 40/121 \\
Q1 & pixel & 121 & 0.562 [0.500, 0.899] & 0/121 \\
\bottomrule
\end{tabular}
\end{table}

\paragraph{Original native checkpoint: box-height adjustment}
\label{app:native_boxheight}
This post-hoc diagnostic uses the frozen full-frame predictions of the
original 1,400-item, v1-trained CourtDyn-native checkpoint, with 140 complete
items per clip/family. The covariate $h_i$ is the median raw MOT box height
at exactly the four presented frame times, mapped by the image identifiers
and scaled to the $960\times540$ rendered canvas. It does not use neighbouring
frames or the all-window height ruler. We assign average ranks to ties and
separately regress the prediction ranks and one-decimal v3 reference ranks
on an intercept and the height ranks using ordinary least squares. Their
residual Pearson correlation is the adjusted rank statistic reported in
Section~\ref{sec:gt}. This controls one scalar rank covariate, not all visual
or motion confounders. Observations remain dependent, and Q1 retains the
original checkpoint's training-event overlap; no significance test or
independent-event inference is claimed.

\paragraph{Four-sample annotation-trajectory reference}
\label{app:four_sample_trajectory}
For Q1 and Q2 speed/path, each family retains all 140 questions. We reuse the
original image-plane footpoint average over a $\pm5$-frame neighbourhood
(up to 11 available annotation frames), before homography projection. The dense path joins 13 points at five-frame spacing;
the four-point path uses the original offsets $(0,20,40,60)$, with identical
endpoints. Each speed equals its corresponding path divided by the actual
$2.002$ s. Thus only the polyline sampling changes; the original v3 and
$960\times540$ pixel targets remain the scoring references.
Table~\ref{tab:four_sample_trajectory} reports unrounded four-point/dense
path ratios. For scoring, four-point answers are formatted to one decimal;
we retain the original T-MRA $T$ and $\epsilon$, including their pixel
scaling by $K$. All eight cells score $100$, with correlations of
$0.9978$--$0.9989$ against their original scoring references. The largest
relative path deficit is $33.4\%$: tolerance coverage is not exact recovery.
These are chord deficits along smoothed annotations, which may also contain
annotation noise. Because the smoothing uses neighbouring frames not shown
to the model, this analysis does not test image-coordinate readability.

\begin{table}[t]
\centering\small
\setlength{\tabcolsep}{3pt}
\caption{Four-sample/dense annotated path ratios, median [IQR]. Each row
contains 140 questions per coordinate convention; the $>10\%$ path-loss
counts coincide for ground and pixel coordinates. Speed uses the same path
ratio. Q1 has training-event overlap for the original native checkpoint.}
\label{tab:four_sample_trajectory}
\begin{tabular}{llccc}
\toprule
Clip & Family & Ground ratio [IQR] & Pixel ratio [IQR] & Loss $>10\%$ \\
\midrule
Q2 & speed & $0.990$ [$0.966, 0.998$] & $0.989$ [$0.965, 0.998$] & 7/140 \\
Q2 & path & $0.991$ [$0.960, 0.998$] & $0.991$ [$0.960, 0.998$] & 8/140 \\
Q1 & speed & $0.979$ [$0.951, 0.997$] & $0.979$ [$0.951, 0.997$] & 14/140 \\
Q1 & path & $0.986$ [$0.952, 0.998$] & $0.986$ [$0.952, 0.998$] & 13/140 \\
\bottomrule
\end{tabular}
\end{table}

\section{Benchmark construction and reference details}
\label{supp:benchmark}

\subsection{Families}

CourtDyn is generated from multi-object tracking trajectories of basketball
games~\cite{teamtrack2024}. It contains seven 30-second clips---three from
overhead cameras and four from side cameras---drawn from four quarters of a
single venue, yielding \textbf{4{,}705 questions}. Each question presents four
frames sampled over a 2.0-second window. The player being asked about is drawn
with a red box; a second player, used only by the comparison families, is drawn
with a blue box; the ball is circled in yellow.

The main pool contains three numeric families and two categorical families.
Numeric targets integrate smoothed tracking positions over each window, rather
than only the three displacements visible between the four supplied frames.
Table~\ref{tab:families} gives the numeric scoring parameters. The categorical
visual-pair evaluation uses separately generated, filtered pairs; its retained
counts can differ from the main-pool target of 70 pairs per clip.

\begin{table}[t]
\centering
\caption{CourtDyn families. Numeric $T$ and $\epsilon$ parametrise T-MRA (main text)
  and carry the answer unit; categorical accuracy uses neither parameter.
  Counts describe the main pool, not every filtered visual-pair pool.}
\label{tab:families}
\small
\begin{tabular}{lllll}
\toprule
Family & Main-pool count/clip & Unit & $T$ & $\epsilon$ \\
\midrule
speed & 140 & m/s & 0.30 & 1.0 \\
path & 140 & m & 0.50 & 2.0 \\
timing & 102--140 & s & 0.50 & 2.0 \\
accel & 70 pairs & categorical & --- & --- \\
faster & 70 pairs & categorical & --- & --- \\
\bottomrule
\end{tabular}
\end{table}

\subsection{Ground truth, and which ruler defines it}
\label{sec:gt}

Converting a trajectory in pixels into metres requires a scale. We compare a box-height heuristic with a calibrated ground-plane estimate;
they are not equally valid physical scales in overhead views.

\textbf{v1, the height ruler.} The prompt itself supplies a cue: ``assume a
typical player on this court is 1.93\,m tall''. Taking metres per pixel as
$1.93$ divided by the bounding-box height defines the v1 heuristic. The height cue does not explicitly instruct this
calculation, and a bounding-box height need not be a projected vertical body
height. This heuristic is also used by the metric in the static
setting~\cite{courtsi2026}.

\textbf{v3, the court homography.} Fitting a ground-plane homography to the
court's painted lines gives a ground-plane scale independent of a player's assumed height,
subject to calibration and foot-point errors.

On \emph{side} views the two agree to within about 15\%. On \emph{overhead}
views they differ by roughly a factor of two (median $v3/v1$ ratio
$0.443$--$0.556$), for a simple reason: seen from above, a player's bounding box
is not their height. The rank correlation between the two ground truths is
$0.79$--$0.99$, which indicates substantial but imperfect agreement, not identical rankings.
The original CourtDyn-native first-frame analysis uses v3; historical and
supplemental first-frame comparisons retain v1. Ruler and unit comparisons use
v3 and pixel ground truth, respectively. We identify the scoring convention for each result.

The homography builder estimates court corners automatically from sampled
frames using court-colour masks, white-line fits and cross-frame aggregation.
We manually checked the resulting homography on all seven clips. Its
\emph{shape} is correct---the three-point line, the free-throw circle and the
free-throw lane all register against three independent anchors---but the four
corner correspondences were taken from the outer edge of the painted court
rather than from the white line, estimated to produce a scale inflation of about $1\%$ (upper bound $2.4\%$).
A strictly uniform within-clip rescaling would leave rank correlations unchanged;
the visual check alone does not establish that all calibration errors are uniform.

\textbf{The timing family has no usable v3.} The ball is not on the ground, so
projecting its centre through a ground-plane homography misplaces it by metres
(median $v3/v1 = 2.3$). We report timing under v1 only and exclude it from the
main conclusions.

\textbf{Box-height leakage.} The homography does not remove the correlation
between ground truth and box height: it is $+0.42$ on the original side clip,
and its sign flips across clips ($-0.21$ to $+0.42$), so the relationship is clip-dependent and cannot be removed by a single
global correction. We report every result per clip and never pool clips. For the original
CourtDyn-native checkpoint, a post-hoc adjustment for median box height at
the four presented frames changes Q2 speed/path correlations from
$0.764$/$0.760$ to $0.760$/$0.761$; Q1 changes from $0.706$/$0.680$ to
$0.703$/$0.656$ (140 items per cell; \ref{app:native_boxheight}). This
single-covariate rank adjustment does not control all confounding or identify
a motion mechanism. The historical source-SFT/GRPO summaries are archived with the released records.

\subsection{Question-pool quality}

Because the pool is generated automatically, we audited it by hand. A
stratified sample of 42 items (seed 42, three per numeric family (speed and path) per clip) was rendered
twice---the full frame the model receives, and a $6\times$ crop centred on the
annotated box, without which a player only a dozen pixels tall cannot be
identified---and judged on three questions: is the red box the player the
question asks about and the same person in all four frames; do the four frames
cover the motion asked about; is the item free of occlusion, out-of-frame and
mis-boxing. Each question passed $41/42$ ($97.6\%$, Wilson interval
$[87.7, 99.6]$).

The single failure mode, formalised as ``a nearly static player inside a
cluster'', accounts for $10.9\%$ of the corpus and is about six times more
frequent on side views ($12.9$--$19.6\%$) than overhead ($1.1$--$3.6\%$).
Removing those items and recomputing changes $\rho$ by a median of $-0.005$
overhead and $+0.018$ on side views---inconsistent in direction and far smaller
than the roughly $0.5$ gap between the two viewpoints. This sensitivity analysis does not explain away the viewpoint difference;
it does not exclude other item-quality effects or certify the unaudited families.

% ===========================================================================
\subsection{Vision-side minimal pairs}

The usual minimal pair swaps the two options in the text of a single image,
which blocks only half of the degenerate strategies. We instead pair two clips
that share no frames, whose question text is byte-identical and whose answer is
flipped, and score the pair jointly. A model that always names the same referent
scores zero by construction.

\subsection{An image-based oracle on the original side clip}
\label{sec:e0}

We test an image-based reference on Q4\_side\_480--510. At a nominal
$597\times336$ resolution approximating the configured pixel budget, a colour
threshold locates the overlaid red box; its height and the prompt's height cue
provide the v1 scale. This image size is hard-coded in the oracle, not a logged
processor output for every model. A later CPU replay of the local
Qwen3.5-4B processor with the configured pixel cap produced a $20\times36$
patch grid with patch size 16, corresponding to $576\times320$ pixels.
Historical inference grids were not saved; the replay is not a historical
tensor record. The published oracle values below have not been recomputed
at the replayed dimensions.

Against v1 answers, the oracle attains T-MRA $99.8$ and $\rho=0.99$ on speed,
$99.6$ and $0.99$ on path, $98.3$ and $0.99$ on timing, $100.0$ paired on faster
and $89.7$ paired on the filtered visual accel pairs, with $100\%$ box
detection. These results establish useful information in the annotated pixels
for this clip and convention. They do not establish equivalent accuracy under
v3 on all seven clips or all interventions. Four-frame sampling also differs
from the denser tracking trajectory used for ground truth. Because the oracle
reads the artificial overlay, it does not by itself distinguish understanding
of player motion from reading box motion.

\paragraph{Annotation-based sampling check}
A separate CPU check replaces the dense annotated path by a polyline through
the same four source sampling times (\ref{app:four_sample_trajectory}). In
Q2, the median retained path fraction is $0.9895$--$0.9908$ across speed/path
and ground/pixel coordinates. All eight Q1/Q2 cells reach T-MRA $100$ under
the original tolerances despite nonzero path deficits, reaching $33.4\%$ in
one item. High tolerance scores therefore do not imply exact path recovery.
The sampled coordinates use the original smoothing over neighbouring
annotation frames, including frames hidden from the model; this is not an
image-based oracle or evidence that the model can read those coordinates.

\section{Additional analyses}
\label{supp:additional}

\subsection{Post-hoc cross-question consistency}

Speed and path questions share 100 matched trajectory windows in Q2. On
this identical cohort, the median ratio $\hat d/(\Delta t\hat v)$ changes
from $0.999$ [IQR $0.999$, $1.155$] under metre prompts to
$0.500$ [$0.500$, $0.659$] under pixel prompts. A common multiplicative
coordinate rescaling cancels in this ratio, although question-specific
conventions remain possible. Counts compatible with one-decimal rounding
in each requested unit are $42/100$ and $4/100$, respectively
(\ref{app:consistency}). These tolerances represent different physical
precisions, so the counts are not a unit-invariant comparison.
This descriptive check adds no independent events and does not identify an
internal mechanism; constant answers can also be self-consistent.

\subsection{Matched training-label results}
\label{sec:matched_label_results}

Table~\ref{tab:matched_label_main} compares both new adapters against the
same exact v3 reference. Both held-out events belong to one game. The
v3-trained adapter has higher metre rank correlations in all four cells,
but T-MRA decreases in three: Q1 speed by $2.21$ points, and Q2
speed/path by $4.43$/$10.07$; Q1 path increases by $0.29$.
The v1-trained adapter exceeds the test-label constant oracle by only
$2.43$--$2.93$ points in three cells and $14.07$ in Q2 path. For the
v3-trained adapter, Q1 speed exceeds it by $0.36$ points and Q2 speed
falls below it by $2.00$. These are descriptive margins, not significant
improvements; the oracle uses test labels and is not a deployable baseline.
Rescoring identical full predictions against exact v3 instead of v1 raises
T-MRA by $23.64$--$46.71$ points and lowers mean absolute error (MAE) in
all eight cells. This is reference sensitivity, not a model improvement.
Across models, v3 training gives higher MAE in all four cells under
\emph{both} references (\ref{app:paired_reference_scores}).
These are seed-42 differences. The three-seed sweep of the main text shows which
of them survive seed variation: v3 training gives higher path $\rho$ and lower
speed T-MRA than v1 training in every seed, whereas the v1 path T-MRA alone spans
$60.5$--$74.2$ across seeds, so the Q2 path difference of $10.07$ points is not a
label effect.

\begin{table}[t]
\centering\footnotesize
\setlength{\tabcolsep}{3pt}
\caption{Matched adapters, full frames and retained height cue. v1/v3 name
the training labels; all metre scores use exact v3 references. Each cell has
140 observed and parsed answers. $M/O$ is model T-MRA / best test-label
constant T-MRA; $R$ is the median paired cm/m prediction ratio (expected
$100$), with 140 positive metre denominators.}
\label{tab:matched_label_main}
\begin{tabular}{lrrrrrr}
\toprule
 & \multicolumn{3}{c}{speed} & \multicolumn{3}{c}{path} \\
\cmidrule(lr){2-4}\cmidrule(lr){5-7}
Train/event & $M/O$ & $\rho$ & $R$ & $M/O$ & $\rho$ & $R$ \\
\midrule
v1 / Q1 & 54.14/51.57 & 0.600 & 1.50 & 47.21/44.29 & 0.512 & 1.60 \\
v3 / Q1 & 51.93/51.57 & 0.663 & 1.38 & 47.50/44.29 & 0.778 & 1.00 \\
v1 / Q2 & 72.07/69.64 & 0.715 & 1.56 & 74.21/60.14 & 0.694 & 1.14 \\
v3 / Q2 & 67.64/69.64 & 0.758 & 1.57 & 64.14/60.14 & 0.773 & 1.00 \\
\bottomrule
\end{tabular}
\end{table}

The paired cm/m medians remain $1.00$--$1.60$, versus the exact factor
$100$. Projected-pixel/metre medians are $1.00$--$3.47$, whereas median
reference ratios are $27.06$--$27.50$; projected and ground motion remain
different quantities (\ref{app:matched_label_results}). On the same
12 text-only probes, v1 training gives $9/12$ correct (cm $3/6$, mm $6/6$),
and v3 training gives $4/12$ (cm $1/6$, mm $3/6$), compared with the original
checkpoint's $2/12$. All text answers parse. On the $246$ diagnostic items of
the expanded probe the order is the same, v1 above v3 above the original
checkpoint, and all three sit significantly below the unadapted backbone
(main text, text-probe table); both new adapters also apply a supplied pixel
factor correctly to a stated number while their image pixel/metre ratios stay
far below the required $27.5$.

First-frame repetition lowers all eight metre T-MRA scores by
$17.93$--$41.07$ points. The seven defined common-cohort correlations
also decrease (\ref{app:matched_label_results}). v3 Q1 speed has
135/140 valid control outputs and a correlation of $0.076$; v3 Q2 speed
has 137/140, all equal to $0.1$, so its control correlation and correlation
difference are undefined. The eight invalid answers receive zero T-MRA.
All other visual cells parse completely. This establishes sensitivity to
the tested temporal input change within two events, not a specific motion
mechanism or cross-game replication.

\subsection{Ordering and first-frame sensitivity by viewpoint}
\label{sec:e1a}

Table~\ref{tab:viewpoint} reports \rhosp{} per clip, with the arrow pointing to
the first-frame$\times 4$ control. Figure~\ref{fig:viewpoint} shows the same
data.

\begin{table}[t]
\centering
\caption{Spearman \rhosp{} against v1 per clip, full $\rightarrow$ first-frame$\times 4$.
  Reported per clip and never pooled. Each arrow pairs full and control on that clip.}
\label{tab:viewpoint}
\footnotesize
\setlength{\tabcolsep}{4pt}
\begin{tabular}{llllll}
\toprule
Clip & SFT speed & SFT path & GRPO speed & GRPO path \\
\midrule
Q1\_top\_0-30      & $0.59 \to 0.01$ & $0.80 \to -0.13$ & $0.57 \to 0.04$ & $0.77 \to 0.05$ \\
Q2\_top\_480-510   & $0.51 \to 0.05$ & $0.66 \to -0.17$ & $0.58 \to -0.03$ & $0.62 \to -0.10$ \\
Q4\_top\_0-30      & $0.45 \to -0.06$ & $0.65 \to 0.07$ & $0.49 \to -0.06$ & $0.66 \to -0.01$ \\
\midrule
Q1\_side\_0-30     & $0.23 \to 0.14$ & $0.51 \to -0.12$ & $0.30 \to 0.03$ & $0.51 \to -0.10$ \\
Q4\_side\_570-600  & $0.34 \to 0.15$ & $0.41 \to 0.19$ & $0.33 \to 0.18$ & $0.34 \to 0.02$ \\
Q4\_side\_480-510  & $0.09 \to 0.28$ & $-0.00 \to 0.17$ & $0.16 \to 0.31$ & $0.02 \to 0.30$ \\
Q2\_side\_300-330  & $-0.25 \to -0.35$ & $-0.12 \to -0.19$ & $-0.22 \to -0.26$ & $-0.07 \to -0.30$ \\
\bottomrule
\end{tabular}
\end{table}

\begin{figure}[t]
\centering
\includegraphics[width=\linewidth]{fig1_viewpoint_rho.pdf}
\caption{Rank correlation against v1 per clip and model, overhead and side panels. Filled
  markers are the full four frames, open markers the first-frame$\times 4$
  control, joined per cell. Plotted from recorded outputs and stated references using AI-assisted Python code (see the AI-use declaration of the main text).}
\label{fig:viewpoint}
\end{figure}

On the three overhead clips, predictions show positive rank agreement:
\rhosp{} is $0.45$--$0.80$ and decreases under the first-frame control, with
$\Delta\rho = 0.46$--$0.93$ for source-SFT and $0.54$--$0.73$ for GRPO, and
itemwise bootstrap intervals on the control that contain zero. Controlling for
median box height leaves the partial correlation at $0.45$--$0.78$, which does not support box height alone as an explanation. The three camera
clips repeat the pattern within one venue; they are not three independent venues.

The side views are not a uniform failure. The regularity is sharper than
``motion is not read from the side'': \emph{clips whose full-frame \rhosp{} is
already above zero lose it when motion is removed, and clips whose \rhosp{} is
already near zero or negative have nothing to lose.} Q1\_side and
Q4\_side\_570 give $\Delta\rho$ of $0.09$--$0.63$ and $0.15$--$0.32$;
Q2\_side is flat; on the original side clip the control scores higher.
This prediction was written down before the side-view controls were run, and it
held.


\noindent The remaining analyses of the earlier version of this study (additional checkpoints and cameras, the choice of reference ruler, scale and unit comparisons across checkpoints, the zoom intervention, paired families, reproduction arms and a second basketball venue) are archived verbatim with the released code (\texttt{legacy\_analyses\_v1\_S4.tex} in the repository) and are not needed for the claims of the main text.

\section{Per-backbone, read-then-convert, probe and zero-shot results}
\label{supp:revision}

\subsection{Per-backbone results at seed 42}
\label{supp:backbones42}

The verdicts of these paragraphs are those of the main text's backbone table;
the numbers are the per-cell values behind them.

Qwen2.5-VL-3B met the criterion at epoch four. Its metre-only adapter now
ranks motion, with $\rho$ from $0.26$ to $0.51$, and converts in no cell
($R=0.63$--$1.00$). R1 is intermediate because only the two Q2 cells clear
$\rho\ge0.30$; the Q1 cells reach $0.26$ and $0.295$. R4 is untestable: with the
first frame repeated, both speed cells collapse to a single answer, so $\rho$
is undefined, while in the two path cells $\rho$ falls from $0.295$ to $0.08$ and
from $0.51$ to $0.07$. The mixed-unit adapter returns $R=100$ in every
non-degenerate cell, but its centimetre answers track the reference in only one
of three explicit cells ($\rho_{\mathrm{cm}}=0.45$) and two of four in-template
cells ($0.33$, $0.48$). Under the amended rule, then, R2 is intermediate: the
factor of 100 is there, but reading is not shown. R3 replicates: metre-only
training lowers diagnostic accuracy from $56.1\%$ to $28.5\%$
($p=1.7\times10^{-15}$), and mixed-unit training leaves it at $52.4\%$
($p=0.26$).

InternVL3-2B, run under the round-2 protocol with its own unadapted-backbone
cells, met the criterion at epoch two. It reproduces the ordering and its
dependence on motion: metre $\rho$ is $0.33$--$0.71$, and first-frame
repetition lowers it in all four cells ($\Delta\rho=0.26$--$0.42$), so R4
replicates. Its path answers do not follow the unit ($R=1.10$ and $1.20$); its
speed answers move a little ($R=2.20$ and $3.67$). One speed cell has a
degenerate centimetre arm and the other exceeds $R\le3$, so R1 is
intermediate. The mixed-unit adapter converts path length with centimetre
answers that track the reference ($R=92$--$100$, $\rho_{\mathrm{cm}}=0.45$--$0.73$
in both wordings). Its centimetre speed answers take two or three distinct
values, as on Qwen3.5-4B; R2 therefore replicates in the training wording, where
the Q2 speed cell has three values and $\rho_{\mathrm{cm}}=0.71$, and is untestable
with explicit prompts. The arithmetic result does not transfer: metre-only
training raises diagnostic accuracy from $64.2\%$ to $72.0\%$ ($p=0.0094$), and
mixed-unit training lowers it to $50.8\%$ ($p=0.00061$), so R3 fails.

Pixtral-12B, from an unrelated model family, met the criterion at epoch two.
Its metre-only adapter orders motion in all four cells ($\rho=0.36$--$0.63$) and
converts in none ($R=1.00$--$1.05$), so R1 replicates. First-frame repetition
lowers $\rho$ in both speed cells ($\Delta\rho=0.35$ and $0.27$); in both path
cells the repeated frame collapses the answers to a single value, so $\rho$ is
undefined and R4 is untestable. The mixed-unit adapter gives ratios of $100$ and
$143$, but its centimetre answers track the reference in two of three
non-degenerate explicit cells ($\rho_{\mathrm{cm}}=0.36$ and $0.55$, both path),
so R2 is intermediate with explicit prompts and untestable in the training
wording. R3 replicates: metre-only training lowers diagnostic accuracy from
$49.2\%$ to $31.7\%$ ($p=4.3\times10^{-6}$), and mixed-unit training raises it to
$62.2\%$ ($p=3.1\times10^{-4}$).

Idefics3-8B met the criterion at epoch two on the development clip, which
points to scale rather than architecture for the degeneracy of SmolVLM2-2.2B at
the recipe's cap. On the test clips, however, its metre-only adapter gives two
distinct metre answers in three of four cells, so R1 and R4 are untestable. Its
mixed-unit adapter gives $R=115$--$169$ in both wordings, but the centimetre
answers do not track the reference ($\rho_{\mathrm{cm}}\le0.27$ in every cell),
so R2 is intermediate: the factor of 100 without reading, the failure the
round-2 amendment was written for. Text arithmetic falls from $58.5\%$ to
$42.3\%$ under metre-only training ($p=9.0\times10^{-8}$) and to $44.3\%$ under
mixed-unit training ($p=1.5\times10^{-4}$), so R3 is intermediate. The
degeneracy does not recur with the seed: at seed 43 all four metre-only cells
are live (3--8 distinct metre answers, $\rho=0.22$--$0.43$) and none converts
($R\le1.13$), so R1 replicates; first-frame repetition lowers $\rho$ in all four
cells and text arithmetic falls from $58.5\%$ to $43.5\%$
($p=7.5\times10^{-7}$) under metre-only training, so R4 and R3 replicate as
well. At seed 44 the cells are live again but R1 and R4 are intermediate, so
across the three seeds R1 and R4 are not established and R3 holds at two of
three (three-seed table of the main text).

Gemma3-12B, also selected at epoch two, runs against the unit half of the
claim at seed 42. Its metre-only adapter orders motion in all four cells
($\rho=0.31$--$0.51$), and first-frame repetition lowers $\rho$ in all four
($\Delta\rho=0.21$--$0.31$), so R4 replicates. But asked for centimetres it
converts: $R=12.21$ and $98.94$ on Q1, $2.02$ and $13.91$ on Q2, three of four
cells at or above $R\ge10$, so R1 fails. The conversion is partial (the
centimetre answers track the reference with $\rho_{\mathrm{cm}}=0.21$--$0.33$,
below the metre arm), but it is there without any centimetre supervision.
Its text arithmetic is not damaged either:
$89.4\%$ unadapted, $92.3\%$ after metre-only training ($p=0.039$) and $88.6\%$
after mixed-unit training ($p=0.77$), so R3 fails. Its mixed-unit adapter
repeats the speed pattern of Qwen3.5-4B: the centimetre speed answers take two
distinct values in both explicit cells (R2-E untestable), while path converts
($R=92$ and $100$). The conversion does not recur with the seed: at seed 43
only the two path cells convert, and at seed 44 no cell does ($R\le6.79$), so
R1 is intermediate and then replicates, and its label across the three seeds
is seed-dependent (three-seed table of the main text); the arithmetic stays undamaged at every
seed.

\subsection{Read-then-convert}
\label{supp:rtc}

The read-then-convert (RTC) adapter is trained on the mixed-unit pool with
every target rewritten as ``\emph{reading in metres} m $\rightarrow$
\emph{value unit}'', so that one answer carries the metre reading and the
converted value. The decisive test is conversion to units absent from
training: km/h and ft/s for speed, ft and mm for path, on both test clips
(eight cells). A cell converts when the paired ratio of the converted answer
to the metre reading lies within $[f/1.5,\,1.5f]$ of the exact factor $f$,
the arm is non-degenerate and the metre reading keeps $\rho\ge0.30$; the
preregistered hypothesis H1 holds with at least six of eight cells. Two
controls were recorded with it: the metre-only and mixed-unit adapters
prompted in the RTC output format without RTC training (H6, the format
control), and the RTC adapter's metre ranking against the metre-only
adapter's (H4). The replication at seeds 43 and 44 and on InternVL3-2B,
Pixtral-12B and Gemma3-12B was specified before those runs.
Table~\ref{tab:rtc} gives the hypotheses per adapter and the per-cell
ratios; the main text summarises the outcome as a prompting effect of
Qwen3.5-4B.

% Generated by tools/make_paper2_extra_tables.py -- do not edit by hand.
\begin{table}[t]
\centering\footnotesize
\setlength{\tabcolsep}{3pt}
\caption{Read-then-convert: preregistered hypotheses per adapter.}
\label{tab:rtc}
\begin{tabular}{lcccccc}
\toprule
Adapter & H1 & H3 & H4 & H5 & H6 v3 & H6 M1 \\
\midrule
Qwen3.5-4B, seed 42 & holds (8) & fixed (2) & fails & untestable (2) & holds (6) & holds (6) \\
Qwen3.5-4B, seed 43 & intermediate (5) & fixed (2) & holds & holds (4) & holds (6) & intermediate (5) \\
Qwen3.5-4B, seed 44 & intermediate (4) & fixed (2) & fails & holds (4) & holds (6) & holds (6) \\
InternVL3-2B & fails (0) & partial (1) & fails & holds (4) & fails (0) & fails (0) \\
Pixtral-12B & fails (1) & partial (1) & fails & untestable (2) & fails (0) & intermediate (5) \\
Gemma3-12B & intermediate (3) & partial (1) & holds & holds (4) & fails (0) & fails (0) \\
\bottomrule
\end{tabular}
\end{table}

\begin{table}[t]
\centering\footnotesize
\setlength{\tabcolsep}{3pt}
\caption{RTC-trained adapters: median answer ratio over the exact factor ($R/f$; 1.00 = exact conversion) per unseen-unit cell, with the metre-arm rank correlation $\rho$ of the same cell in parentheses. U1 = km/h (speed) or ft (path); U2 = ft/s or mm; px = pixels with the ruler sentence. A cell converts under the rule only if $R/f\in[1/1.5,1.5]$, the arm is non-degenerate and $\rho\ge0.30$.}
\label{tab:rtccells}
\resizebox{\linewidth}{!}{%
\begin{tabular}{llcccccc}
\toprule
Adapter & clip & speed U1 & path U1 & speed U2 & path U2 & speed px & path px \\
\midrule
Qwen3.5-4B, seed 42 & Q1 & 0.99 (0.38)* & 1.31 (0.64)* & 1.17 (0.41)* & 1.00 (0.58)* & 0.66 (0.44) & 1.00 (0.65)* \\
Qwen3.5-4B, seed 42 & Q2 & 1.02 (0.46)* & 1.31 (0.64)* & 1.32 (0.56)* & 1.00 (0.64)* & 0.64 (0.01) & 0.98 (0.57)* \\
Qwen3.5-4B, seed 43 & Q1 & 0.28 (0.38) & 1.09 (0.69)* & 1.00 (0.25) & 0.91 (0.72)* & 0.04 (0.49) & 0.81 (0.77)* \\
Qwen3.5-4B, seed 43 & Q2 & 0.28 (0.37) & 1.09 (0.74)* & 1.00 (0.42)* & 0.92 (0.70)* & 0.04 (0.66) & 0.04 (0.53) \\
Qwen3.5-4B, seed 44 & Q1 & 0.28 (0.44) & 1.00 (0.67)* & 1.57 (0.12) & 0.94 (0.69)* & 0.04 (0.57) & 0.03 (0.71) \\
Qwen3.5-4B, seed 44 & Q2 & 0.32 (0.47) & 1.00 (0.72)* & 1.57 (0.37) & 1.00 (0.68)* & 0.03 (0.61) & 0.04 (0.69) \\
InternVL3-2B & Q1 & 0.25 (0.06) & 0.30 (0.26) & 0.30 (0.28) & 0.10 (0.22) & 0.04 (0.25) & 3.57 (0.16) \\
InternVL3-2B & Q2 & 0.28 (-0.14) & 0.30 (0.49) & 0.30 (0.55) & 0.10 (0.39) & 0.04 (0.21) & 3.52 (0.15) \\
Pixtral-12B & Q1 & 1.07 (0.10) & 1.00 (0.29) & 1.00 (0.15) & 0.10 (--) & 1.41 (0.06) & 1.03 (-0.07) \\
Pixtral-12B & Q2 & 1.11 (0.35)* & 1.00 (0.27) & 1.00 (0.27) & 0.10 (--) & 1.24 (0.26) & 1.01 (--) \\
Gemma3-12B & Q1 & 0.28 (0.15) & 1.00 (0.28) & 0.72 (0.10) & 1.00 (0.32)* & 1.00 (0.10) & 1.00 (0.24) \\
Gemma3-12B & Q2 & 0.28 (0.36) & 1.00 (0.24) & 0.72 (0.42)* & 1.00 (0.31)* & 0.98 (0.45)* & 0.98 (0.35)* \\
\bottomrule
\end{tabular}}
\par\vspace{3pt}\begin{flushleft}\footnotesize H1: unseen-unit cells converted (of 8; holds $\ge6$); H3: centimetre speed (of 2); H4: metre $\rho$ not inferior to the metre-only adapter; H5: first-frame repetition lowers $\rho$ (of 4); H6: metre-only (v3) and mixed-unit (M1) adapters in the RTC format without training (of 8).\end{flushleft}
\end{table}
% caption trimmed by tools/trim_table_captions.py


\subsection{Representation probe: verdict stability and nested re-analysis}
\label{supp:probe}

The probe asks whether the reference quantity can be read linearly from the
hidden states, and whether that readout depends on the motion in the frames.
For each of the unadapted backbone, CourtDyn-native, the metre-only v3 and
the mixed-unit adapter, one forward pass over the explicit metre prompt of
all 280 speed and path items of five clips yields hidden states at layers 4,
8, \ldots, 32 under two poolings (mean over the image tokens, and the last
prompt token). A ridge regression on the standardised features predicts the
log-metric target, with the penalty chosen by inner grouped cross-validation
(groups: clip and 2-second window start) and the layer--pooling pair selected
once per model by cross-validated $R^2$ on the overhead clips. H1
(decodability) holds when the grouped cross-validated Spearman correlation on
the overhead clips is at least $0.30$ and exceeds the 99th percentile of 200
within-clip label permutations. H3 (motion dependence) applies the probe fit
on full-frame items to the first-frame$\times4$ features of the Q2 clip and
holds when the correlation drops by at least $0.30$. H2 asks whether the
image-plane or the floor-plane quantity is the better-explained target. The
preregistered analysis used one cross-validation fold seed; because H2
depended on the fold assignment, the analysis was repeated post hoc over ten
fold seeds, and Table~\ref{tab:probe} reports how often each verdict was
reached.

The preregistered pipeline has a selection step outside the cross-validation:
the layer--pooling pair is chosen on all overhead items before H1 is scored on
the same items, the permutation null keeps that chosen feature and penalty,
and the first-frame items of H3 are scored by a probe fit on all items,
including their own four-frame versions. A post hoc re-analysis
(\texttt{tools/repr\_probe\_nested.py}) moves the selection and fitting
steps inside an outer five-fold grouped cross-validation: within each training
fold the feature is selected by inner cross-validated $R^2$ and the penalty by
inner cross-validation, the 200 within-clip label permutations repeat this
whole procedure, and the first-frame$\times4$ items of each held-out fold are
scored by that fold's probe, so full and static predictions come from
identical training exposure. The one quantity still computed on all overhead
items is the per-feature standardisation (column means and standard
deviations of the hidden states), which uses no labels.
Table~\ref{tab:probe_nested} sets the two pipelines side
by side. H1 holds in every model and family under both. H3 does not: under
the preregistered pipeline the unadapted backbone and the mixed-unit adapter
kept their correlation under first-frame repetition ($\rho_{\mathrm{static}}$
$0.52$--$0.63$), whereas under the nested pipeline the correlation drops by
$0.40$--$0.92$ in all four models. The contrast between models does not
survive the corrected procedure, and the main text no longer relies on it.

% Generated by tools/make_paper2_extra_tables.py -- do not edit by hand.
\begin{table}[t]
\centering\footnotesize
\caption{Representation probe: verdict frequency over 10 fold seeds.}
\label{tab:probe}
\begin{tabular}{lcccccc}
\toprule
Model & H1 speed & H1 path & H3 speed & H3 path & H2 path & H2 speed \\
\midrule
unadapted & 10/10 & 10/10 & 0/10 & 0/10 & UNTESTABLE 7, floor-plane 3 & UNTESTABLE 10 \\
CourtDyn-native & 10/10 & 10/10 & 9/10 & 10/10 & excluded 10 & excluded 10 \\
metre-only v3 & 10/10 & 10/10 & 10/10 & 10/10 & image-plane 7, UNTESTABLE 2, floor-plane 1 & UNTESTABLE 10 \\
mixed-unit & 10/10 & 10/10 & 0/10 & 0/10 & UNTESTABLE 8, floor-plane 2 & UNTESTABLE 8, image-plane 2 \\
\bottomrule
\end{tabular}
\par\vspace{3pt}\begin{flushleft}\footnotesize Counts over 10 cross-validation fold seeds. H1: the linear readout beats the permutation null; H3: the readout loses its correlation under first-frame repetition; H2: image-plane vs.\ floor-plane quantity (untestable when the control does not separate them).\end{flushleft}
\end{table}
% caption trimmed by tools/trim_table_captions.py


% Generated by tools/repr_probe_nested_table.py -- do not edit by hand.
\begin{table}[t]
\centering\footnotesize
\setlength{\tabcolsep}{4pt}
\caption{Representation probe: preregistered pipeline and the nested re-analysis.}
\label{tab:probe_nested}
\begin{tabular}{llcccccc}
\toprule
 & & \multicolumn{3}{c}{preregistered (selection on all items)} & \multicolumn{3}{c}{nested (selection inside folds)} \\
\cmidrule(lr){3-5}\cmidrule(lr){6-8}
Model & family & $\rho_{\mathrm{full}}$ & $\rho_{\mathrm{static}}$ & dep. & $\rho_{\mathrm{full}}$ & $\rho_{\mathrm{static}}$ & dep. \\
\midrule
unadapted & speed & 0.71 & 0.56 & no & 0.67 & 0.15 & yes \\
unadapted & path & 0.80 & 0.63 & no & 0.75 & 0.36 & yes \\
CourtDyn-native & speed & 0.76 & 0.00 & yes & 0.74 & 0.07 & yes \\
CourtDyn-native & path & 0.77 & 0.09 & yes & 0.78 & 0.11 & yes \\
metre-only v3 & speed & 0.77 & -0.01 & yes & 0.76 & 0.09 & yes \\
metre-only v3 & path & 0.80 & -0.03 & yes & 0.75 & -0.17 & yes \\
mixed-unit & speed & 0.76 & 0.52 & no & 0.75 & 0.17 & yes \\
mixed-unit & path & 0.79 & 0.63 & no & 0.76 & 0.02 & yes \\
\bottomrule
\end{tabular}
\par\vspace{3pt}\begin{flushleft}\footnotesize Q2\_top, 140 items per family. $\rho_{\mathrm{full}}$: out-of-fold Spearman of the probe against log metre on the four-frame items; $\rho_{\mathrm{static}}$: the same probe applied to the first-frame$\times4$ features of the same items; dep.\ = motion-dependent ($\rho_{\mathrm{static}}\le\rho_{\mathrm{full}}-0.30$). Preregistered: feature chosen on all overhead items, static items scored by a probe fit on all items. Nested: feature and penalty chosen inside each outer training fold, static items scored by the probe of the fold that holds them out; H1 holds for every row against 200 full-pipeline permutations.\end{flushleft}
\end{table}


\subsection{Zero-shot sweep of public backbones}
\label{supp:t36}

Nine public backbones (Qwen3.5-4B in NF4 and bf16, Qwen2.5-VL-3B and -7B,
InternVL3-2B and -8B, Gemma3-4B, LLaVA-OneVision-7B and SmolVLM2-2.2B) answer
the explicit metre, centimetre and pixel cells of both test clips without
any fine-tuning, under the same prompts, image policies and greedy decoding
as the adapters, plus the first-frame$\times4$ control on the metre arm and
the 12-item text conversion probe. Table~\ref{tab:t36} reports, per
backbone and cell, the number of distinct parsed metre answers, the metre
rank correlation, the paired centimetre/metre and pixel/metre answer ratios,
and whether both the metre and centimetre arms are live. No decision rule
was recorded for this sweep; it is descriptive. One cell is invalid: the bf16
Qwen3.5-4B run was loaded through a generic path that left the model's
thinking mode on, so no answer was produced within the token budget; the
cell is marked rather than rerun.

% Generated by tools/make_paper2_extra_tables.py -- do not edit by hand.
\begin{table}[t]
\centering\footnotesize
\setlength{\tabcolsep}{3pt}
\caption{Zero-shot sweep of public backbones on the explicit unit cells.}
\label{tab:t36}
\resizebox{\linewidth}{!}{%
\begin{tabular}{llcccccc}
\toprule
Model & clip & family & distinct & $\rho$ & $R_{\mathrm{cm}}$ & $R_{\mathrm{px}}$ ($K$) & live \\
\midrule
gemma3-4b & Q1 & speed & 8 & 0.03 & 0.3 & 2.9 (27) & no \\
gemma3-4b & Q1 & path & 5 & -0.17 & 1.2 & 12.1 (27) & yes \\
gemma3-4b & Q2 & speed & 9 & -0.23 & 0.3 & 2.9 (27) & no \\
gemma3-4b & Q2 & path & 7 & -0.23 & 1.0 & 10.7 (27) & yes \\
internvl3-2b & Q1 & speed & 3 & -0.05 & 8.5 & 10.0 (27) & yes \\
internvl3-2b & Q1 & path & 2 & 0.33 & 1.0 & 10.2 (27) & no \\
internvl3-2b & Q2 & speed & 2 & 0.02 & 8.5 & 10.0 (27) & no \\
internvl3-2b & Q2 & path & 2 & 0.20 & 6.4 & 10.2 (27) & no \\
internvl3-8b & Q1 & speed & 9 & 0.15 & 19.0 & 16.7 (27) & yes \\
internvl3-8b & Q1 & path & 11 & 0.50 & 58.3 & 47.2 (27) & yes \\
internvl3-8b & Q2 & speed & 8 & 0.10 & 19.0 & 16.7 (27) & yes \\
internvl3-8b & Q2 & path & 9 & 0.18 & 35.0 & 37.9 (27) & yes \\
llava-ov-7b & Q1 & speed & 4 & 0.03 & 0.8 & 0.5 (27) & yes \\
llava-ov-7b & Q1 & path & 3 & -0.03 & 52.6 & 17.4 (27) & no \\
llava-ov-7b & Q2 & speed & 4 & -0.11 & 1.0 & 1.6 (27) & yes \\
llava-ov-7b & Q2 & path & 3 & -0.04 & 52.6 & 17.4 (27) & no \\
qwen25vl-3b & Q1 & speed & 3 & -0.02 & 8.2 & 10.4 (27) & no \\
qwen25vl-3b & Q1 & path & 2 & -0.01 & 10.0 & 10.0 (27) & no \\
qwen25vl-3b & Q2 & speed & 3 & 0.13 & 9.9 & 10.4 (27) & no \\
qwen25vl-3b & Q2 & path & 2 & 0.10 & 10.0 & 10.0 (27) & no \\
qwen25vl-7b & Q1 & speed & 3 & 0.26 & 9.9 & 2.4 (27) & no \\
qwen25vl-7b & Q1 & path & 4 & 0.29 & 104.2 & 10.4 (27) & yes \\
qwen25vl-7b & Q2 & speed & 3 & 0.35 & 25.4 & 2.4 (27) & yes \\
qwen25vl-7b & Q2 & path & 4 & 0.15 & 104.2 & 10.4 (28) & yes \\
qwen35-4b & Q1 & speed & 1 & -- & 1.0 & 3.0 (27) & no \\
qwen35-4b & Q1 & path & 3 & 0.25 & 0.0 & 0.0 (27) & no \\
qwen35-4b & Q2 & speed & 1 & -- & 1.0 & 3.0 (27) & no \\
qwen35-4b & Q2 & path & 3 & 0.32 & 0.0 & 0.0 (27) & no \\
qwen35-4b-bf16 & Q1 & speed & 0 & -- & -- & -- (--) & INVALID \\
qwen35-4b-bf16 & Q1 & path & 0 & -- & -- & -- (--) & INVALID \\
qwen35-4b-bf16 & Q2 & speed & 0 & -- & -- & -- (--) & INVALID \\
qwen35-4b-bf16 & Q2 & path & 0 & -- & -- & -- (--) & INVALID \\
smolvlm2-2b & Q1 & speed & 5 & 0.22 & 1.6 & 1.0 (27) & yes \\
smolvlm2-2b & Q1 & path & 4 & -0.22 & 7.3 & 731.4 (27) & yes \\
smolvlm2-2b & Q2 & speed & 4 & -0.00 & 1.2 & 1.0 (27) & yes \\
smolvlm2-2b & Q2 & path & 2 & -0.24 & 5.7 & 1015.8 (27) & no \\
\bottomrule
\end{tabular}}
\par\vspace{3pt}\begin{flushleft}\footnotesize distinct: parsed metre answers; $\rho$: metre rank correlation; $R_{\mathrm{cm}}$, $R_{\mathrm{px}}$: median paired ratios (exact 100 and the clip's $K\approx27$); live: both arms $\ge3$ distinct answers. INVALID: harness failure (thinking mode left on), not a model result.\end{flushleft}
\end{table}
% caption trimmed by tools/trim_table_captions.py


\subsection{Per-item conversion of the mixed-unit adapters}
\label{supp:item_ratios}

The preregistered unit-response rule judges a cell by the median paired
ratio $R$. A median near $100$ is also compatible with two separately
calibrated output mappings, so Table~\ref{tab:item_ratios} reports, for each
seed, test clip and family, the distribution of the item-level ratio
$r_i = \hat y_{\mathrm{cm},i}/(100\,\hat y_{\mathrm{m},i})$ and the share of
items within a factor of two and within a factor of $1.25$ of exact
conversion (post hoc, from the recorded predictions). For path, $r_i$ lies
within a factor of two in 99--100\% of items in all six cells; the share within
a factor of $1.25$ is 69--89\% at seeds 42 and 44 and 19--27\% at seed 43,
whose median ratio is $76$. For speed the centimetre answers take one to three
values, so $r_i$ is within a factor of $1.25$ of exact conversion in at most
1\% of items.

% Generated by tools/mixunit_item_ratios.py -- do not edit by hand.
\begin{table}[t]
\centering\footnotesize
\setlength{\tabcolsep}{4pt}
\caption{Per-item centimetre/metre answer ratios of the mixed-unit adapters, explicit prompts.}
\label{tab:item_ratios}
\begin{tabular}{llcccccc}
\toprule
Seed & clip & family & distinct cm & median $r_i$ & IQR & within $\times2$ & within $\times1.25$ \\
\midrule
42 & Q1 & speed & 2 & 1.83 & [1.83, 2.28] & 74\% & 0\% \\
42 & Q1 & path & 8 & 1.00 & [1.00, 1.14] & 100\% & 86\% \\
42 & Q2 & speed & 2 & 2.50 & [1.83, 2.50] & 44\% & 0\% \\
42 & Q2 & path & 7 & 1.00 & [1.00, 1.09] & 100\% & 89\% \\
43 & Q1 & speed & 1 & 1.83 & [1.57, 1.83] & 99\% & 1\% \\
43 & Q1 & path & 6 & 0.76 & [0.69, 0.76] & 100\% & 19\% \\
43 & Q2 & speed & 2 & 1.83 & [1.83, 1.83] & 86\% & 0\% \\
43 & Q2 & path & 6 & 0.76 & [0.69, 1.00] & 99\% & 27\% \\
44 & Q1 & speed & 1 & 1.83 & [1.83, 1.83] & 100\% & 1\% \\
44 & Q1 & path & 5 & 1.00 & [0.88, 1.17] & 100\% & 76\% \\
44 & Q2 & speed & 3 & 1.83 & [1.83, 1.83] & 96\% & 0\% \\
44 & Q2 & path & 5 & 1.00 & [1.00, 1.27] & 100\% & 69\% \\
\bottomrule
\end{tabular}
\par\vspace{3pt}\begin{flushleft}\footnotesize $r_i = \hat y_{\mathrm{cm},i}/(100\,\hat y_{\mathrm{m},i})$ over the 140 items of each cell with both answers parsed; exact conversion gives $r_i=1$. ``within $\times2$'': $0.5\le r_i\le2$; ``within $\times1.25$'': $0.8\le r_i\le1.25$. The preregistered rule uses the median ratio $R$ of the same pairs.\end{flushleft}
\end{table}


% ===========================================================================

\bibliographystyle{elsarticle-num}
\bibliography{refs}
\end{document}
